\documentclass[11pt]{article}
\usepackage{fontspec}
\usepackage[french,english]{babel}
\usepackage{amsmath}
\usepackage{amssymb}
\usepackage{amsthm}
\usepackage{mathrsfs}
\usepackage{tikz-cd}
\usetikzlibrary{babel}
\usepackage{enumitem}

\newcommand{\printedpage}[1]{\marginpar{\footnotesize p.~#1}}
\newcommand{\msp}[1]{\marginpar{\footnotesize [#1]}}
\newcommand{\barlongleftrightarrow}{\mathrel{\text{\ooalign{$\longleftrightarrow$\cr\hidewidth$\mid$\hidewidth\cr}}}}
\newcommand{\idx}[2]{\noindent #1\dotfill #2\par}

\begin{document}

\printedpage{49}
% head: A. Grothendieck : Dérivateurs, Ch. XIII, Édition des Universités de Montpellier II et Paris VII
% foot: Version 25 janvier 2004

\msp{page 91}Let
\[ \mathcal{C} \;=\; \underline{M}^\pi/\mathcal{X} \]
be the following category. Its objects are the triples
\[ (Z,i,\alpha) \quad \text{such that} \quad Z\in M^\pi,\ i\in I,\ \alpha : Z\longrightarrow X_i \ , \]
the arrows
\[ (Z,i,\alpha) \;\longrightarrow\; (Z',i',\alpha') \]
are the pairs $(u,\lambda)$ with
\[ u : Z\longrightarrow Z', \qquad \lambda : i\longrightarrow i', \qquad (\textit{i.e.}\ \text{we have } i\le i' \text{ and } \lambda \text{ is the unique arrow from } i \text{ to } i'), \]
such that the diagram
\[
\begin{tikzcd}[row sep=large, column sep=large]
Z \arrow[r,"\alpha"] \arrow[d,"u"'] & X_i \arrow[d,"\text{transition } \lambda_{i',i}"] \\
Z' \arrow[r,"\alpha'"] & X_{i'}
\end{tikzcd}
\]
commutes. We can therefore define $\mathcal{C}$ as the fiber product in the cartesian square
\[
\begin{tikzcd}[row sep=large, column sep=large]
 & \mathcal{C} \arrow[r,"\text{cofibered}"] \arrow[d] \arrow[dl,dotted,"\mathcal{Z}"'] & I \arrow[d,"\mathcal{X} : X \longmapsto X_i"] \\ % [sic] X \longmapsto X_i -> i \longmapsto X_i, the functor \mathcal{X} goes from I, so the variable should be i
\underline{M}^\pi & \underline{M}^\pi/E \arrow[l,dotted,"\text{source}"'] \arrow[r,"\text{target}","\text{cofibered}"'] & E.
\end{tikzcd}
\]
The fiber $\mathcal{C}_i$ of $\mathcal{C}$ at $i\in I$ is $\mathrm{Filt}^\pi/X_i$ \texttt{[rather} $M^\pi/X_i$\texttt{]}. We have, as in \textit{b}) $\Longrightarrow$ first half of \textit{a}), that $\underline{M}^\pi$ is a
\msp{page 92}system of strict generators of $\underline{M}$ \texttt{[this is not clear without further hypotheses on} $\underline{M}$\texttt{]}, and $X_i\in\operatorname{Ob}\underline{M}$ by hypothesis,
\[ X_i \;=\; \varinjlim_{Z\in\operatorname{Ob}\underline{M}^\pi/X_i} Z \;=\; \varinjlim_{\mathcal{C}_i}\mathcal{Z}|\mathcal{C}_i, \]
so the given direct system $(X_i)_{i\in I}$ on $I$ with values in $X$, is deduced from the direct system (not ordered in general) on $\mathcal{C}$, with values in $\underline{M}^\pi$, by ``integrating over the fibers'' of $\mathcal{C}\longrightarrow I$. Besides, since $\underline{M}^\pi/E\longrightarrow E$ is cofibered, hence proper, the same is true of $\mathcal{C}\longrightarrow I$, so the $\varinjlim$ of $\mathcal{Z}$ over the fibers $\mathcal{C}_i$ ($=$ $\underline{M}^\pi/X_i$) do form a covariant functor $I\longrightarrow E$, whose direct limit is that of $\mathcal{Z}$ (\footnote[43]{\textbf{NB} If we admit the validity of $a$), then it already follows that $\varinjlim\limits_{i\in I} X_i \simeq \varinjlim\limits_{\alpha\in C} Z_\alpha$ is in $M$. (Stability of $M$ under $\varinjlim$ large relative to $\pi$).}).

\printedpage{50}
% head: A. Grothendieck : Dérivateurs, Ch. XIII, Édition des Universités de Montpellier II et Paris VII
% foot: Version 25 janvier 2004

\msp{page 93}Let, for every $i\in I$, $\tilde{\mathcal{C}}_i$ be the ordered set
\[ \tilde{\mathcal{C}}_i = \{\text{subcategories } \mathcal{C}'_i \text{ of } \mathcal{C}_i = \underline{M}^\pi/X_i \mid \operatorname{card}\mathcal{Z}(\mathcal{C}'_i)\le\pi,\ \ \mathcal{C}'_i \text{ filtering}\}. \]
I admit for the moment the

\textbf{Lemma} (\footnote[44]{Uses the tacit hypothesis that the elements of $\mathrm{Filt}^\pi$ are $\pi$-accessible.}). \textit{For every $X\in M$, the category $\underline{M}^\pi/X$ is filtering. (Under hypothesis} b), \textit{or that every object of $M$ is of the form stated in} a).\textit{)}

It follows, since $X_i\in M$, that $\mathcal{C}_i$ is filtering, so the union of the $\mathcal{C}'_i\in\tilde{\mathcal{C}}_i$ is $\mathcal{C}_i$; the ordered set $\tilde{\mathcal{C}}_i$ is filtering, and every filtering subset $(\mathcal{C}'_{i_\alpha})_{\alpha\in A}$ of cardinality $\le\pi$ of $\tilde{\mathcal{C}}_i$ has a supremum $\mathcal{C}'_{i_A}$ there, namely the union of the $\mathcal{C}'_{i_\alpha}$. $\tilde{\mathcal{C}}_i$ is therefore a $\pi$-adapted ordered set.

We can regard the ordered sets $\tilde{\mathcal{C}}_i$, for varying $i\in I$, as the fibers of an increasing map of ordered sets
\[ p \;:\; \tilde{\mathcal{C}} \;\longrightarrow\; \mathcal{C}, \]
by writing
\[ \left\{ \begin{array}{l}
(\mathcal{C}'_i\subset\mathcal{C}_i) \;\le\; (\mathcal{C}'_j\subset\mathcal{C}_j) \\
\text{if and only if } i\le j, \\
\text{and the canonical transition functor } \mathcal{C}_i\longrightarrow\mathcal{C}_j \text{ maps } \mathcal{C}'_i \text{ into } \mathcal{C}'_j.
\end{array}\right. \]
We thus see that $\tilde{\mathcal{C}}$ too is filtering and large relative to $\pi$.

\msp{page 94}We define a direct system on $\tilde{\mathcal{C}}$, with values in $\underline{\mathrm{Filt}}^\pi$, by setting
\[ Z_{\mathcal{C}'_i} \;=\; \varinjlim_{\alpha\,\in\,\mathcal{C}'_i\ (\subset\,\mathcal{C}_i\,=\,\underline{M}^\pi/X_i)} Z_\alpha \]
(where an $\alpha\in\mathcal{C}_i$ consists, I recall, of $Z_\alpha\xrightarrow{u_\alpha}X_i$, with $Z_\alpha\in M^\pi\subset\underline{\mathrm{Filt}}^\pi$). These $Z_{\mathcal{C}'_i}$ are no longer necessarily in $M^\pi$, but only in $\mathrm{Filt}^\pi$. (Then what was the use of defining the $\mathcal{C}_i$ themselves as $\underline{M}^\pi/X_i$, instead of $\underline{\mathrm{Filt}}^\pi/X_i$?)

The direct system
\[ (Z_{\mathcal{C}'_i})_{\mathcal{C}'_i\in\tilde{\mathcal{C}}_i} \]
is perfectly well adapted to the cardinal filtration for the cardinal $\pi$, and it is immediate that its $\varinjlim$ is isomorphic to $\varinjlim\limits_{\alpha\in\operatorname{Ob}\mathcal{C}_i=\underline{M}^\pi/X_i} Z_\alpha = X_i\in M$ . Thus by \textit{b}), the set $\tilde{\mathcal{C}}_M$ of the $\mathcal{C}'_i\in\tilde{\mathcal{C}}_i$ such that $Z_{\mathcal{C}'_i}\in M$ (so $Z_{\mathcal{C}'_i}\in M^\pi$) is cofinal in $\tilde{\mathcal{C}}_i$. It follows therefore that the set $\tilde{\mathcal{C}}_M$ is also cofinal in $\tilde{\mathcal{C}}$.

\printedpage{51}
% head: A. Grothendieck : Dérivateurs, Ch. XIII, Édition des Universités de Montpellier II et Paris VII
% foot: Version 25 janvier 2004

\msp{page 95}We also check immediately that
\[ \varinjlim_{\alpha\in\tilde{\mathcal{C}}} Z_\alpha \;=\; \varinjlim X_i \;=\; X, \]
so, $\tilde{\mathcal{C}}_M$ being cofinal,
\[ \varinjlim_{\alpha\in\tilde{\mathcal{C}}_M} Z_\alpha \;=\; X. \]
All in all, we have managed to represent $X$ in the form
\[ X \;=\; \varinjlim_{\alpha\in A} X_\alpha, \]
where this time the $X_\alpha$ are \textit{in} $M^\pi$ (not just in $M$), and $A$ is large relative to $\pi$. But we do not know that every filtering subset $B$ of cardinality $\le\pi$ of $A$ has a supremum $\alpha_B$ in $A$, and still less that $X_{\alpha_B} = \varinjlim\limits_{\alpha\in B} X_\alpha$ (\footnote[45]{If we assume $(a_\pi)$, this argument proves at least that $M$ is stable under $\varinjlim$ large relative to $\pi$.}). So we cannot yet apply \textit{b}) directly. If it were only a matter of proving the second half of \textit{a}), where we already assume the $X_i\in M^\pi$ and $I$ large relative to $\pi$, we have gained nothing at all by our labors!

\msp{page 96}But let us return to this case
\[ X \;=\; \varinjlim X_i \;\left|\; \begin{array}{l} X_i\in M^\pi \\[1.5ex] I \text{ large relative to } \pi, \end{array}\right. \]
and introduce again
\[ \tilde{I} \;=\; \{I'\subset I \mid I' \text{ filtering, of cardinality } \le\pi\}, \]
so this time $\tilde{I}$ is indeed adapted to $\pi$ (filtering, and every filtering subset of cardinality $\le\pi$ has a supremum), and if we set
\[ X_{\tilde{\imath}} \;=\; \varinjlim_{i\in I'_{\tilde{\imath}}} X_i, \]
we see that the direct system
\[ (X_{\tilde{\imath}})_{\tilde{\imath}\in\tilde{I}} \]
is perfectly adapted to $\pi$ and the cardinal filtration. Moreover, I say that the set $\tilde{I}_M$ is \textit{cofinal} in $\tilde{I}$. Indeed, $\tilde{I}_M$ contains the subset $\tilde{I}_0$ consisting of the $I'\in\tilde{I}$ which have a greatest element -- and $\tilde{I}_0$ is indeed cofinal! Now, we can apply
\printedpage{52}
% head: A. Grothendieck : Dérivateurs, Ch. XIII, Édition des Universités de Montpellier II et Paris VII
% foot: Version 25 janvier 2004
\msp{page 97}the ``it suffices'' in \textit{b}), which concludes that $X = \varinjlim X_i\in M$, q.e.d. % [sic] conclue -> conclut, verb form (indicative of conclure)

We have proved $b\Longrightarrow a$, without using the hypothesis that every object of $\mathrm{Filt}^\pi$ is $\pi$-accessible. But we have used this hypothesis to prove that $b$ (or $a$) implies that $M$ is stable under $\varinjlim$ large relative to $\pi$ (lemma 2, p. 90).

$a\overset{?}{\Longrightarrow} b$. The ``it suffices'' in \textit{b}) is an immediate consequence of \textit{a}). It remains to prove that if, in the $\pi$-adapted direct system $(X_i)_{i\in I}$ with limit $X$, we have $X\in M$, then $I_M$ is cofinal in $I$. But by hypothesis \textit{a}), we will have $X = \varinjlim\limits_J Y_j$, with $J$ large relative to $\pi$, the $Y_j$ in $M^\pi$. Let us admit that $\mathrm{Filt}^\pi\subset E_\pi$, \textit{i.e.} every element of $\mathrm{Filt}^\pi$, in particular the $X_i$, are $\pi$-accessible.

I think I will not manage to prove $a_\pi\Longrightarrow b_\pi$, but only this: If $\pi'$ is a cardinal $>\pi$ (say the successor of $\pi$ in the ordinal sequence of cardinals), then $a_\pi\Longrightarrow b_\pi$.

\msp{page 98}This will follow from

\textbf{Lemma 3.} \textit{Let $X\simeq\varinjlim\limits_I X_i$, $X\simeq\varinjlim\limits_J Y_j$ be two representations of $X$ in the form of filtering $\varinjlim$, with $I$ large relative to $\pi'$, $J$ large relative to $\pi$, with $\pi$, $\pi'$ two cardinals such that $\pi'>\pi$. We assume $(X_i)_{i\in I}$ $\pi'$-adapted,} i.e. \textit{for every filtering subset $I'$ of $I$ of cardinality $\le\pi'$, $I'$ has a supremum $i_{I'}$, and $X_{i_{I'}} = \varinjlim\limits_{i\in I'} X_i$. Then for every $i\in I$, we can find a subset $I'$ of $I$, containing $i$,} large relative to $\pi$, \textit{of cardinality $\le\pi'$, and a subset $J'$ of $J$, finally an isomorphism $I'\xrightarrow[\sim]{\lambda}J'$ and a strictly increasing map $\mu : J'\longrightarrow I'$ and homomorphisms} \texttt{[end of chapter]}

\newpage
\printedpage{53}
% head: A. Grothendieck : Dérivateurs, Ch. XIII, Édition des Universités de Montpellier II et Paris VII
% foot: Version 25 janvier 2004

\section*{Index of notations}

\idx{LLP, RLP}{1}
\idx{$\Phi_*$, $\Psi^*$}{1}
\idx{$\widetilde{\Phi}$, $\overline{\Psi}$}{1}
\idx{$\Phi \longleftrightarrow \Psi$}{3}
\idx{$\mathrm{Cof}_W$, $\mathrm{Fib}_W$}{8}
\idx{$W^{\mathrm{univ}}$}{9}
\idx{$W_{\mathrm{univ}}$}{10}
\idx{$\underline{\mathrm{Fib}}_W S\subset\mathcal{M}/S$}{10}
\idx{$W_S$}{10}
\idx{$W_S^{\mathrm{u}}$}{10}
\idx{$S\backslash\underline{\mathrm{Cof}}_W$}{10}
\idx{$\Phi \barlongleftrightarrow \Psi$}{28}
\idx{$\mathcal{M}_{\mathrm{cf}}$, $W_{\mathrm{cf}}$}{29}
\idx{$\mathrm{Ho}_W$}{29}
\idx{$\mathrm{Hom}_W$}{30}
\idx{$\pi_{\mathrm{d}}$, $\pi$}{30}
\idx{$\mathcal{M}_{\mathrm{c}}$}{30}
\idx{$W_{\mathrm{c}}$}{31}
\idx{$\pi_{\mathrm{g}}$}{31}
\idx{$\mathcal{M}_{\mathrm{f}}$}{31}
\idx{$f\sim g$}{32}
\idx{$\Psi,\ \Phi,\ \mathcal{M}(\Psi),\ W(\Psi),\ \mathcal{M}(\Phi),\ W(\Phi)$}{33}
\idx{$\mathcal{M}(I)$, $W(I)$}{34}
\idx{$F(I),\ TF(I),\ C(I),\ TC(I)$}{35}
\idx{$\alpha_s$}{35}
\idx{$u^*$, $u_*$, $u_!$}{45}
\idx{$u^*_W$, $u^W_*$, $u^W_!$}{49}
\idx{$\mathrm{HO}_W(I)$}{49}
\idx{Diag}{50}
\idx{$\mathrm{HO}_{\mathcal{M},W}(I)$, $\mathbf{D}_{\mathcal{M},W}(I)$}{54}
\idx{$\mathrm{HO}^{\mathrm{lc}}_W(I)$, $u^*_{\mathrm{lc}}$}{56}
\idx{$u^{\mathrm{f}}_*$}{59}
\idx{$u^*_{\mathrm{c}}$}{67}
\idx{$\mathfrak{U}$}{77}
\idx{$\mathrm{Filt}^\pi(E)$}{77}

\end{document}
